% ============================================================================
% MÓDULO 2 - Manipulación de datos con dplyr y tidyr
% ============================================================================
\chapter{Manipulación de datos con dplyr y tidyr}
\label{cap:manipulacion}

\begin{objetivos}
\begin{itemize}
  \item Explicar qué son los \emph{datos ordenados} (\emph{tidy data}) y por
        qué facilitan todo el análisis posterior.
  \item Importar archivos CSV con \paquete{readr}, controlar los tipos de
        columna y detectar problemas de lectura.
  \item Encadenar pasos con el pipe (\codigo{\%>\%} y \codigo{|>}) y leer tu
        código como si fuera una receta.
  \item Responder preguntas de negocio con los verbos de \paquete{dplyr}:
        \codigo{filter()}, \codigo{select()}, \codigo{mutate()},
        \codigo{arrange()}, \codigo{summarise()} y \codigo{group\_by()}.
  \item Calcular KPIs con funciones de ventana: crecimiento mes contra mes,
        acumulado del año, rankings y promedios móviles.
  \item Manejar fechas con \paquete{lubridate} y texto con \paquete{stringr}.
  \item Unir tablas con los distintos tipos de \emph{join} y detectar sus
        trampas (llaves duplicadas, registros huérfanos).
  \item Cambiar la forma de una tabla con \paquete{tidyr} (formato ancho y
        largo, tablas dinámicas, celdas vacías).
  \item Limpiar un archivo ``sucio'' paso a paso y construir un reporte
        mensual completo que se exporta a Excel.
\end{itemize}
\end{objetivos}

Es lunes, 9:00 de la mañana. La directora comercial te escribe: \emph{``¿Me
mandas antes de la junta de las 12 las ventas de 2023 por mes y por
categoría, cuánto crecimos contra el mes anterior, quiénes son nuestros cinco
mejores clientes y qué tienda va peor? En Excel, por favor.''} Tienes los
datos, pero están repartidos en varios archivos: las ventas por un lado, el
catálogo de productos por otro, los clientes, las tiendas y los empleados en
archivos distintos. Ninguno, por sí solo, responde la pregunta.

Si haces esto a mano en Excel, te espera una mañana de \codigo{BUSCARV},
tablas dinámicas, copiar y pegar\ldots{} y el próximo mes, repetirlo todo. En
este módulo aprenderás a hacerlo con \paquete{dplyr} y \paquete{tidyr}, los
dos paquetes que más vas a usar en tu vida como analista. Al final del módulo
(sección~\ref{sec:m2-caso}) construirás exactamente ese reporte con un script
de R que, el mes que entra, se vuelve a ejecutar en segundos.

Este es el módulo más largo del libro, y con razón: aquí está el 80\% de tu
trabajo diario. No intentes memorizar todas las funciones; concéntrate en
entender la \emph{lógica} de cada verbo y teclea los ejemplos tú mismo. Todo
el código usa los datos del curso (\archivo{ventas\_retail.csv},
\archivo{productos.csv}, \archivo{clientes.csv}, \archivo{tiendas.csv},
\archivo{empleados.csv} y \archivo{transacciones.csv}), así que puedes
reproducir cada salida exactamente como aparece en el libro.

% ----------------------------------------------------------------------------
\section{Preparar los datos: el 80\% del trabajo}
\label{sec:m2-motivacion}

En la industria se repite mucho una cifra: un analista dedica alrededor del
80\% de su tiempo a \emph{preparar} datos (conseguirlos, limpiarlos, unirlos y
darles forma) y solo el 20\% restante a analizarlos y presentarlos. La cifra
exacta varía según la empresa, pero la idea es cierta: los datos casi nunca
llegan listos para el análisis. La figura~\ref{fig:m2-flujo} muestra el flujo
de trabajo típico y qué paquete del \emph{tidyverse} te ayuda en cada paso.

\begin{figure}[H]
\centering
\begin{tikzpicture}[
  paso/.style={draw=primario, fill=primario!8, rounded corners=3pt,
               minimum width=2.25cm, minimum height=1cm, align=center,
               font=\sffamily\small\bfseries, text=primario},
  paq/.style={font=\ttfamily\scriptsize, text=gristexto, align=center},
  flecha/.style={-{Stealth[length=2.5mm]}, thick, color=acento}]
  \node[paso] (imp) {Importar};
  \node[paso, right=0.55cm of imp] (ord) {Ordenar};
  \node[paso, right=0.55cm of ord] (tra) {Transformar};
  \node[paso, right=0.55cm of tra] (res) {Resumir};
  \node[paso, right=0.55cm of res, draw=exito, fill=exito!8, text=exito]
       (com) {Comunicar};
  \foreach \a/\b in {imp/ord, ord/tra, tra/res, res/com}
    \draw[flecha] (\a) -- (\b);
  \node[paq, below=0.2cm of imp] {readr\\readxl};
  \node[paq, below=0.2cm of ord] {tidyr};
  \node[paq, below=0.2cm of tra] {dplyr\\lubridate\\stringr};
  \node[paq, below=0.2cm of res] {dplyr\\(group\_by,\\summarise)};
  \node[paq, below=0.2cm of com] {ggplot2\\writexl\\rmarkdown};
  \begin{scope}[on background layer]
    \node[draw=acento, dashed, rounded corners=5pt, inner sep=6pt,
          fit=(imp)(res), label={[font=\sffamily\scriptsize\bfseries,
          text=acento]above:Este módulo: preparar los datos}] {};
  \end{scope}
\end{tikzpicture}
\caption{Flujo de trabajo de un análisis de datos y los paquetes que usarás
en cada paso. Este módulo cubre los cuatro primeros; el Módulo~3 se dedica a
comunicar con gráficas.}
\label{fig:m2-flujo}
\end{figure}

\subsection{Datos ordenados (\emph{tidy data})}
\label{sec:m2-tidy}

Antes de transformar datos conviene saber a qué forma queremos llegar. El
estadístico Hadley Wickham, creador de buena parte de los paquetes de este
módulo, propuso una forma estándar llamada \termino{datos ordenados}
(\emph{tidy data}).

\begin{concepto}{Las tres reglas de los datos ordenados}
Una tabla está \emph{ordenada} cuando cumple tres reglas:
\begin{enumerate}
  \item \textbf{Cada variable es una columna.} Fecha, sucursal, producto y
        monto son cuatro variables: cuatro columnas.
  \item \textbf{Cada observación es una fila.} Una venta, un cliente o un
        mes de una sucursal ocupa una fila.
  \item \textbf{Cada valor está en su propia celda.} Nada de
        ``\texttt{CDMX-Centro}'' en una celda si ciudad y zona son cosas
        distintas.
\end{enumerate}
Cuando tus datos cumplen estas reglas, todas las funciones del
\emph{tidyverse} ``encajan'': filtras, agrupas, unes y graficas sin pelearte
con la forma de la tabla.
\end{concepto}

El ejemplo más común de datos \emph{desordenados} en una empresa es el reporte
de Excel ``bonito'', pensado para leerse, no para analizarse. En la
tabla~\ref{tab:m2-tidy} el reporte de la izquierda tiene los meses como
columnas: la variable \emph{mes} está escondida en los encabezados y la
variable \emph{ventas} está repartida en tres columnas. La versión ordenada
(derecha) tiene una columna por variable.

\begin{table}[H]
\centering
\caption{El mismo dato en dos formas: un reporte ``ancho'' típico de Excel
(desordenado) y su versión ordenada (\emph{tidy}).}
\label{tab:m2-tidy}
\begin{minipage}[t]{0.44\textwidth}
\centering
\textbf{\sffamily\small\color{peligro}Desordenado (ancho)}\\[4pt]
\small
\begin{tabular}{lrrr}
\toprule
Sucursal & Ene & Feb & Mar \\
\midrule
Centro & 12\,500 & 13\,200 & 11\,800 \\
Norte  &  9\,800 & 10\,100 & 11\,350 \\
Sur    & 15\,200 & 14\,900 & 16\,050 \\
\bottomrule
\end{tabular}

\smallskip
{\footnotesize\color{gristexto} 3 filas; la variable ``mes''\\ vive en los
encabezados.}
\end{minipage}%
\hfill
\begin{minipage}[t]{0.48\textwidth}
\centering
\textbf{\sffamily\small\color{exito}Ordenado (largo)}\\[4pt]
\small
\begin{tabular}{llr}
\toprule
sucursal & mes & ventas \\
\midrule
Centro & Ene & 12\,500 \\
Centro & Feb & 13\,200 \\
Centro & Mar & 11\,800 \\
Norte  & Ene &  9\,800 \\
\multicolumn{3}{c}{\ldots} \\
Sur    & Mar & 16\,050 \\
\bottomrule
\end{tabular}

\smallskip
{\footnotesize\color{gristexto} 9 filas; cada fila es una sucursal en un mes.}
\end{minipage}
\end{table}

¿Por qué es mejor la versión larga? Porque preguntas como ``¿cuánto se vendió
en total en febrero?'' o ``¿qué mes fue el mejor de cada sucursal?'' se
resuelven con un filtro o un agrupamiento, sin importar si mañana el reporte
trae 3 meses o 36. En la sección~\ref{sec:m2-tidyr} aprenderás a pasar de una
forma a la otra con una sola función. Otros síntomas de datos desordenados que
encontrarás en la práctica son:

\begin{itemize}
  \item Dos datos en una celda: \texttt{"Guadalajara, Jal."} o
        \texttt{"2023-Q1"}.
  \item Celdas combinadas de Excel: la región aparece solo en la primera fila
        del bloque y las demás quedan vacías.
  \item Totales y subtotales mezclados con los datos (la fila ``Total
        general'' al final).
  \item Varias tablas en una misma hoja, o encabezados en dos renglones.
\end{itemize}

\subsection{El \emph{tidyverse}: una caja de herramientas coherente}

El \termino{tidyverse} es una colección de paquetes de R diseñados para
trabajar juntos con datos ordenados. Comparten la misma filosofía: funciones
con nombres de verbos, el \emph{data frame} siempre como primer argumento y
resultados que también son \emph{data frames}. La tabla~\ref{tab:m2-tidyverse}
resume los que usarás en este módulo.

\begin{table}[H]
\centering
\caption{Paquetes del \emph{tidyverse} que usarás en este módulo.}
\label{tab:m2-tidyverse}
\small
\begin{tabularx}{\textwidth}{@{}lXl@{}}
\toprule
Paquete & ¿Para qué sirve? & Funciones estrella \\
\midrule
\paquete{readr}     & Leer y escribir archivos de texto (CSV)
                    & \codigo{read\_csv()}, \codigo{write\_csv()} \\
\paquete{tibble}    & La versión moderna del \emph{data frame}
                    & \codigo{tibble()}, \codigo{glimpse()} \\
\paquete{dplyr}     & Transformar: filtrar, calcular, resumir, unir
                    & \codigo{filter()}, \codigo{mutate()}, \codigo{left\_join()} \\
\paquete{tidyr}     & Cambiar la forma de las tablas y tratar vacíos
                    & \codigo{pivot\_longer()}, \codigo{fill()} \\
\paquete{lubridate} & Fechas y horas
                    & \codigo{ymd()}, \codigo{month()}, \codigo{floor\_date()} \\
\paquete{stringr}   & Texto (cadenas de caracteres)
                    & \codigo{str\_detect()}, \codigo{str\_squish()} \\
\paquete{forcats}   & Variables categóricas (factores)
                    & \codigo{fct\_reorder()}, \codigo{fct\_lump()} \\
\paquete{ggplot2}   & Gráficas (Módulo~3)
                    & \codigo{ggplot()}, \codigo{geom\_col()} \\
\paquete{purrr}     & Repetir una operación sobre muchos elementos
                    & \codigo{map()} \\
\bottomrule
\end{tabularx}
\end{table}

Si todavía no lo tienes, instala el \emph{tidyverse} y el paquete
\paquete{writexl} (que usaremos para exportar a Excel). Solo se hace una vez
por computadora:

\begin{rnoejecutar}
# Se ejecuta UNA sola vez (tarda unos minutos)
install.packages(c("tidyverse", "writexl"))
\end{rnoejecutar}

Y al inicio de cada script cargas los paquetes (esto sí, cada vez):

\begin{rcode}
# Carga dplyr, tidyr, readr, stringr, lubridate, ggplot2, etc.
library(tidyverse)
# writexl no es parte del tidyverse: se carga aparte
library(writexl)
\end{rcode}

Al cargar el \emph{tidyverse} verás un mensaje de bienvenida con dos partes.
La primera, \emph{Attaching core tidyverse packages}, lista los nueve paquetes
que se cargaron y sus versiones. La segunda, \emph{Conflicts}, avisa que
algunas funciones ``tapan'' (\emph{mask}) a otras con el mismo nombre; por
ejemplo, \codigo{dplyr::filter() masks stats::filter()} significa que, a
partir de ahora, cuando escribas \codigo{filter()} R usará la versión de
\paquete{dplyr} y no la del paquete base \paquete{stats}. No es un error:
es un aviso, y es justo lo que queremos. (En la sección~\ref{sec:m2-errores}
verás qué pasa cuando olvidas cargar el paquete.)

\begin{consejo}
La notación \codigo{paquete::funcion()} llama a una función de un paquete
específico sin ambigüedad, aunque haya otra con el mismo nombre. Por ejemplo,
\codigo{dplyr::filter()} siempre es el filtro de \paquete{dplyr}. La verás
en foros y en código profesional; también sirve para usar una función de un
paquete instalado sin cargarlo con \codigo{library()}.
\end{consejo}

% ----------------------------------------------------------------------------
\section{Importar datos con \texorpdfstring{\paquete{readr}}{readr}}
\label{sec:m2-readr}

\subsection{Tu primera lectura con \texorpdfstring{\codigo{read\_csv()}}{read\_csv()}}

Un \termino{CSV} (\emph{comma-separated values}) es un archivo de texto donde
cada renglón es una fila y las columnas se separan con comas. Es el formato
más común para intercambiar datos entre sistemas: casi cualquier ERP, CRM o
base de datos puede exportar a CSV. La función \codigo{read\_csv()} de
\paquete{readr} lo lee y devuelve un \emph{tibble}.

Recuerda que todas las rutas del libro son relativas a la carpeta del curso
(la que contiene \archivo{datasets/}). Si en RStudio abriste el proyecto del
curso, ya estás ahí; si no, revisa el Módulo~\ref{cap:fundamentos}.

\begin{rcode}
# Leer el archivo de ventas y guardarlo en el objeto "ventas"
ventas <- read_csv("datasets/ventas_retail.csv")
\end{rcode}
\begin{rsalida}
Rows: 365 Columns: 14
-- Column specification ----------------------------------------------
Delimiter: ","
chr   (3): dia_semana, mes, trimestre
dbl  (10): producto_id, cliente_id, cantidad, precio_unitario, des...
date  (1): fecha

i Use `spec()` to retrieve the full column specification for this data.
i Specify the column types or set `show_col_types = FALSE` to quiet this message.
\end{rsalida}

Este mensaje, la \termino{especificación de columnas}, no es un error: es
\paquete{readr} contándote qué hizo. Léelo de arriba abajo:

\begin{itemize}
  \item \texttt{Rows: 365 Columns: 14}: se leyeron 365 filas (una venta por
        día de 2023) y 14 columnas. Compara siempre estos números con lo que
        esperabas; es tu primera prueba de calidad.
  \item \texttt{Delimiter: ","}: el separador que detectó.
  \item \texttt{chr (3)}: tres columnas de texto (\emph{character}).
  \item \texttt{dbl (10)}: diez columnas numéricas (\emph{double}, números
        con decimales). Los identificadores como \codigo{producto\_id}
        también se leyeron como números.
  \item \texttt{date (1)}: una columna de fecha. \paquete{readr} reconoció
        el formato \texttt{2023-01-01} y lo convirtió automáticamente.
\end{itemize}

\paquete{readr} \emph{adivina} el tipo de cada columna mirando las primeras
1000 filas. Casi siempre acierta, pero no siempre (lo verás en un momento).
Cuando ya conoces el archivo, puedes silenciar el mensaje con
\codigo{show\_col\_types = FALSE}, que es lo que haremos en el resto del
libro.

\subsection{Mirar los datos: \texorpdfstring{\emph{tibble}}{tibble} y
\texorpdfstring{\codigo{glimpse()}}{glimpse()}}

Escribe el nombre del objeto para verlo:

\begin{rcode}
ventas
\end{rcode}
\begin{rsalida}
# A tibble: 365 x 14
   fecha      dia_semana producto_id cliente_id cantidad
   <date>     <chr>            <dbl>      <dbl>    <dbl>
 1 2023-01-01 domingo             31         94       10
 2 2023-01-02 lunes               15        169        8
 3 2023-01-03 martes              14        136        3
 4 2023-01-04 miércoles            3        167        2
 5 2023-01-05 jueves              42         81       10
 6 2023-01-06 viernes             50        191        8
 7 2023-01-07 sábado              43         82       10
 8 2023-01-08 domingo             37        106        7
 9 2023-01-09 lunes               14        119        7
10 2023-01-10 martes              25         11        6
# i 355 more rows
# i 9 more variables: precio_unitario <dbl>, descuento_pct <dbl>,
#   tienda_id <dbl>, vendedor_id <dbl>, subtotal <dbl>,
#   descuento_monto <dbl>, total <dbl>, mes <chr>, trimestre <chr>
\end{rsalida}

Un \termino{tibble} es un \emph{data frame} con mejores modales. Fíjate en
lo que hace al imprimirse:

\begin{itemize}
  \item La primera línea dice qué es y su tamaño: \texttt{A tibble: 365 x 14}.
  \item Debajo de cada nombre de columna aparece su tipo entre
        \texttt{< >}: \texttt{<date>}, \texttt{<chr>}, \texttt{<dbl>}.
  \item Muestra solo 10 filas y las columnas que caben en la pantalla; al
        final te dice cuántas filas y columnas quedaron sin mostrar.
  \item Los números se muestran con pocas cifras significativas (por eso
        verás \texttt{3508.} en vez de \texttt{3507.9}). Es solo la
        impresión: el valor guardado conserva todos sus decimales.
\end{itemize}

Un \emph{data frame} clásico, en cambio, imprime \emph{todas} las filas y
parte las columnas en bloques. Con 365 filas eso inunda la consola; por eso
aquí mostramos solo tres:

\begin{rcode}
# Convertimos a data.frame clásico solo para comparar la impresión
as.data.frame(head(ventas, 3))
\end{rcode}
\begin{rsalida}
       fecha dia_semana producto_id cliente_id cantidad
1 2023-01-01    domingo          31         94       10
2 2023-01-02      lunes          15        169        8
3 2023-01-03     martes          14        136        3
  precio_unitario descuento_pct tienda_id vendedor_id subtotal
1          350.79             5         6           5  3507.90
2          282.75             0         5          24  2262.00
3          296.51            10         2          25   889.53
  descuento_monto    total     mes trimestre
1         175.395 3332.505 2023-01        Q1
2           0.000 2262.000 2023-01        Q1
3          88.953  800.577 2023-01        Q1
\end{rsalida}

Para ver \emph{todas} las columnas de un tibble de un vistazo, usa
\codigo{glimpse()}: pone las columnas como renglones, con su tipo y sus
primeros valores. Es la función que más vas a usar para ``conocer'' una
tabla nueva.

\begin{rcode}
glimpse(ventas)
\end{rcode}
\begin{rsalida}
Rows: 365
Columns: 14
$ fecha           <date> 2023-01-01, 2023-01-02, 2023-01-03, 2023-01~
$ dia_semana      <chr> "domingo", "lunes", "martes", "miércoles", "~
$ producto_id     <dbl> 31, 15, 14, 3, 42, 50, 43, 37, 14, 25, 26, 2~
$ cliente_id      <dbl> 94, 169, 136, 167, 81, 191, 82, 106, 119, 11~
$ cantidad        <dbl> 10, 8, 3, 2, 10, 8, 10, 7, 7, 6, 1, 4, 1, 5,~
$ precio_unitario <dbl> 350.79, 282.75, 296.51, 222.48, 218.83, 302.~
$ descuento_pct   <dbl> 5, 0, 10, 0, 0, 0, 15, 0, 0, 15, 0, 15, 0, 0~
$ tienda_id       <dbl> 6, 5, 2, 7, 3, 3, 4, 1, 6, 5, 2, 7, 10, 1, 5~
$ vendedor_id     <dbl> 5, 24, 25, 18, 8, 19, 3, 5, 8, 12, 5, 23, 25~
$ subtotal        <dbl> 3507.90, 2262.00, 889.53, 444.96, 2188.30, 2~
$ descuento_monto <dbl> 175.395, 0.000, 88.953, 0.000, 0.000, 0.000,~
$ total           <dbl> 3332.505, 2262.000, 800.577, 444.960, 2188.3~
$ mes             <chr> "2023-01", "2023-01", "2023-01", "2023-01", ~
$ trimestre       <chr> "Q1", "Q1", "Q1", "Q1", "Q1", "Q1", "Q1", "Q~
\end{rsalida}

Con esto ya sabes qué contiene cada venta: la \codigo{fecha}, qué producto
(\codigo{producto\_id}), a qué cliente, cuántas unidades, a qué precio, con
qué descuento, en qué tienda y qué vendedor la hizo. Las columnas
\codigo{subtotal}, \codigo{descuento\_monto} y \codigo{total} son cálculos
ya hechos, y \codigo{mes} y \codigo{trimestre} son etiquetas de texto.

\begin{consejo}
Si quieres ver más filas de un tibble usa \codigo{print(ventas, n = 30)};
para ver todas las columnas, \codigo{print(ventas, width = Inf)}. En RStudio
también puedes abrir la tabla como hoja de cálculo con \codigo{View(ventas)}
o haciendo clic en su nombre en el panel \menu{Environment}.
\end{consejo}

\subsection{Controlar los tipos con \texorpdfstring{\codigo{col\_types}}{col\_types}}

A veces la adivinanza de \paquete{readr} no es lo que necesitas. Un caso
típico son los identificadores: \codigo{cliente\_id} es una etiqueta, no una
cantidad (sumar identificadores no tiene sentido). Con el argumento
\codigo{col\_types} le dices a \paquete{readr} exactamente cómo leer cada
columna; las que no menciones se siguen adivinando.

\begin{rcode}
# Leer los identificadores como texto y la cantidad como entero
ventas_tipos <- read_csv(
  "datasets/ventas_retail.csv",
  col_types = cols(
    producto_id = col_character(),   # texto
    cliente_id  = col_character(),   # texto
    cantidad    = col_integer(),     # número entero
    .default    = col_guess()        # el resto: que readr adivine
  )
)
# Revisamos solo las columnas que nos interesan
glimpse(select(ventas_tipos, producto_id, cliente_id, cantidad))
\end{rcode}
\begin{rsalida}
Rows: 365
Columns: 3
$ producto_id <chr> "31", "15", "14", "3", "42", "50", "43", "37", "~
$ cliente_id  <chr> "94", "169", "136", "167", "81", "191", "82", "1~
$ cantidad    <int> 10, 8, 3, 2, 10, 8, 10, 7, 7, 6, 1, 4, 1, 5, 10,~
\end{rsalida}

Ahora \codigo{producto\_id} es \texttt{<chr>} y \codigo{cantidad} es
\texttt{<int>} (entero). Las funciones más usadas en \codigo{cols()} son
\codigo{col\_character()}, \codigo{col\_double()}, \codigo{col\_integer()},
\codigo{col\_logical()}, \codigo{col\_date(format = "\%d/\%m/\%Y")} y
\codigo{col\_skip()} (para no leer una columna). Si solo necesitas algunas
columnas de un archivo grande, \codigo{col\_select} lee únicamente esas:
\codigo{read\_csv(ruta, col\_select = c(fecha, total))}.

\begin{advertencia}
Para unir tablas (sección~\ref{sec:m2-joins}), la llave debe tener el
\emph{mismo tipo} en ambas. Si lees \codigo{producto\_id} como texto en
ventas y como número en productos, el \emph{join} fallará con un error
\texttt{Can't join ... due to incompatible types}. En este módulo dejaremos
los identificadores como números en todas las tablas para evitarlo.
\end{advertencia}

\subsection{Cuando algo sale mal: \texorpdfstring{\codigo{problems()}}{problems()}}

Imagina que la sucursal Norte te manda su inventario y alguien escribió
``N/D'' y ``sin dato'' en lugar de un número. Para simular ese archivo, lo
creamos nosotros mismos con \codigo{write\_lines()} (cada elemento es un
renglón del archivo) en una carpeta \archivo{datos/}:

\begin{rcode}
# Crear la carpeta datos/ (si ya existe, no pasa nada)
dir.create("datos", showWarnings = FALSE)
# Escribir un CSV pequeño con dos valores "raros" en existencias
write_lines(c("sucursal,producto,existencias",
              "Centro,P01,120",
              "Norte,P01,N/D",
              "Sur,P02,85",
              "Este,P02,sin dato"),
            "datos/inventario.csv")
# Leerlo exigiendo que existencias sea numérica
inventario <- read_csv("datos/inventario.csv",
                       col_types = cols(existencias = col_double()))
\end{rcode}
\begin{rsalida}
Warning message:
One or more parsing issues, call `problems()` on your data frame for
details, e.g.:
  dat <- vroom(...)
  problems(dat)
\end{rsalida}

\paquete{readr} no se detiene: lee lo que puede, pone \codigo{NA} (valor
faltante) donde no pudo convertir y te avisa. La función \codigo{problems()}
te dice exactamente dónde está cada problema (su columna \codigo{file}
muestra la ruta completa del archivo; aquí la omitimos para ahorrar espacio):

\begin{rcode}
problems(inventario) %>%
  select(row, col, expected, actual)
\end{rcode}
\begin{rsalida}
# A tibble: 2 x 4
    row   col expected actual  
  <int> <int> <chr>    <chr>   
1     3     3 a double N/D     
2     5     3 a double sin dato
\end{rsalida}

Las filas 3 y 5 del archivo (contando el encabezado como fila 1), en la
columna 3, esperaban un número y encontraron \texttt{N/D} y
\texttt{sin dato}. Como sabemos que esos textos significan ``no hay dato'',
la solución limpia es decírselo a \paquete{readr} con el argumento
\codigo{na}:

\begin{rcode}
inventario <- read_csv("datos/inventario.csv",
                       na = c("", "NA", "N/D", "sin dato"),
                       show_col_types = FALSE)
inventario
\end{rcode}
\begin{rsalida}
# A tibble: 4 x 3
  sucursal producto existencias
  <chr>    <chr>          <dbl>
1 Centro   P01              120
2 Norte    P01               NA
3 Sur      P02               85
4 Este     P02               NA
\end{rsalida}

Ya no hay advertencia y \codigo{existencias} es numérica
(\texttt{<dbl>}). Los faltantes son explícitos: nadie los confundirá con
un cero.

\begin{negocio}{CSV exportados desde Excel en español}
Si tu empresa exporta los CSV desde Excel con configuración regional de
México o España, puede que encuentres dos sorpresas. La primera: columnas
separadas por punto y coma (\texttt{;}) y decimales con coma
(\texttt{1234,50}); para eso existe \codigo{read\_csv2()}. La segunda:
acentos que se ven como \texttt{\~{A}\textcopyright} o
\texttt{\~{A}\textpm}, porque el archivo no está
en UTF-8; se corrige indicando la codificación:
\begin{rnoejecutar}
reporte <- read_csv2("datos/reporte_excel.csv")      # ; y decimal ,
clientes_erp <- read_csv("datos/clientes_erp.csv",
                         locale = locale(encoding = "latin1"))
\end{rnoejecutar}
\end{negocio}

\subsection{Cargar todas las tablas del curso}

Ahora sí, carguemos el resto de las tablas que usaremos en el módulo. Todas
se leen igual; solo cambia el nombre del archivo:

\begin{rcode}
productos     <- read_csv("datasets/productos.csv",
                          show_col_types = FALSE)
clientes      <- read_csv("datasets/clientes.csv",
                          show_col_types = FALSE)
tiendas       <- read_csv("datasets/tiendas.csv",
                          show_col_types = FALSE)
empleados     <- read_csv("datasets/empleados.csv",
                          show_col_types = FALSE)
transacciones <- read_csv("datasets/transacciones.csv",
                          show_col_types = FALSE)

# Un pequeño inventario de lo que cargamos
tibble(
  tabla    = c("ventas", "productos", "clientes", "tiendas",
               "empleados", "transacciones"),
  filas    = c(nrow(ventas), nrow(productos), nrow(clientes),
               nrow(tiendas), nrow(empleados), nrow(transacciones)),
  columnas = c(ncol(ventas), ncol(productos), ncol(clientes),
               ncol(tiendas), ncol(empleados), ncol(transacciones))
)
\end{rcode}
\begin{rsalida}
# A tibble: 6 x 3
  tabla         filas columnas
  <chr>         <int>    <int>
1 ventas          365       14
2 productos        50       13
3 clientes        200       10
4 tiendas          10        8
5 empleados       100       12
6 transacciones  1000       14
\end{rsalida}

En BI a esta distinción se le da nombre: \codigo{ventas} y
\codigo{transacciones} son \termino{tablas de hechos} (registran eventos:
cada fila es una venta) y \codigo{productos}, \codigo{clientes},
\codigo{tiendas} y \codigo{empleados} son \termino{tablas de dimensiones}
(describen a los protagonistas de esos eventos). Las conectarás en la
sección~\ref{sec:m2-joins}.

\begin{hazlotu}
Lee \archivo{datasets/transacciones.csv} \emph{sin}
\codigo{show\_col\_types = FALSE} y revisa la especificación de columnas.
¿Qué tipo le asignó \paquete{readr} a la columna \codigo{hora}? Luego usa
\codigo{glimpse()} sobre \codigo{clientes} y responde: ¿de qué tipo es
\codigo{es\_premium}?
\end{hazlotu}

% ----------------------------------------------------------------------------
\section{El pipe: encadenar pasos como una receta}
\label{sec:m2-pipe}

Supón que quieres ver las ventas de la tienda 1 en diciembre, solo con la
fecha, el producto, la cantidad y el total, de la más grande a la más
chica. Son tres pasos: filtrar, seleccionar columnas y ordenar. Hay tres
maneras de escribirlo.

\textbf{Opción 1: funciones anidadas.} Cada función envuelve a la anterior.
Funciona, pero se lee de adentro hacia afuera y los argumentos quedan lejos
de su función:

\begin{rcode}
arrange(
  select(
    filter(ventas, tienda_id == 1, mes == "2023-12"),
    fecha, producto_id, cantidad, total
  ),
  desc(total)
)
\end{rcode}

\textbf{Opción 2: objetos intermedios.} Se lee mejor, pero llenas tu
ambiente de objetos (\codigo{paso1}, \codigo{paso2}\ldots) que no vuelves a
usar, y es fácil equivocarte de nombre:

\begin{rcode}
paso1 <- filter(ventas, tienda_id == 1, mes == "2023-12")
paso2 <- select(paso1, fecha, producto_id, cantidad, total)
paso3 <- arrange(paso2, desc(total))
\end{rcode}

\textbf{Opción 3: el pipe.} El operador \codigo{\%>\%}, llamado
\termino{pipe} (tubería), toma lo que está a su izquierda y lo pasa como
\emph{primer argumento} a la función de su derecha. Así,
\codigo{x \%>\% f(y)} es lo mismo que \codigo{f(x, y)}:

\begin{rcode}
ventas %>%                                   # toma las ventas,
  filter(tienda_id == 1, mes == "2023-12") %>% # y luego filtra,
  select(fecha, producto_id, cantidad, total) %>% # y luego elige,
  arrange(desc(total))                       # y luego ordena
\end{rcode}
\begin{rsalida}
# A tibble: 5 x 4
  fecha      producto_id cantidad  total
  <date>           <dbl>    <dbl>  <dbl>
1 2023-12-18          18        9 2646. 
2 2023-12-26          26        5 1847. 
3 2023-12-12          35        3  441. 
4 2023-12-24          42        1  349. 
5 2023-12-25          32        1   54.0
\end{rsalida}

Las tres opciones dan exactamente el mismo resultado, pero la tercera se lee
de arriba abajo como una receta. Lee cada \codigo{\%>\%} como
\textbf{``y luego''}: \emph{``toma las ventas, y luego filtra la tienda 1 en
diciembre, y luego elige cuatro columnas, y luego ordena por total
descendente''}. La tienda 1 tuvo 5 ventas en diciembre y la mayor fue de
\$2,646 el 18 de diciembre.

\begin{consejo}
El atajo de teclado para escribir el pipe en RStudio es
\tecla{Ctrl} + \tecla{Shift} + \tecla{M} (en Mac, \tecla{Cmd} +
\tecla{Shift} + \tecla{M}). Úsalo desde hoy: lo vas a teclear cientos de
veces al día. Escribe un paso por línea y deja el \codigo{\%>\%} \emph{al
final} de la línea, nunca al principio de la siguiente.
\end{consejo}

\subsection{El pipe nativo \texorpdfstring{\codigo{|>}}{|>}}

Desde la versión 4.1, R trae su propio pipe, el \termino{pipe nativo}
\codigo{|>}, que no necesita cargar ningún paquete. Para el trabajo diario
ambos hacen lo mismo:

\begin{rcode}
# Mismo resultado que con %>%: contar las ventas de diciembre
ventas |>
  filter(mes == "2023-12") |>
  nrow()
\end{rcode}
\begin{rsalida}
[1] 31
\end{rsalida}

Las diferencias que te pueden afectar son pocas:

\begin{itemize}
  \item \codigo{|>} siempre exige paréntesis: \codigo{ventas |> nrow()}
        funciona pero \codigo{ventas |> nrow} da error, mientras que
        \codigo{ventas \%>\% nrow} sí funciona.
  \item Para mandar el dato a un argumento que no es el primero,
        \codigo{\%>\%} usa un punto (\codigo{data = .}) y \codigo{|>} usa un
        guion bajo con nombre (\codigo{data = \_}, desde R 4.2).
  \item \codigo{\%>\%} viene del paquete \paquete{magrittr} y se carga con
        el \emph{tidyverse}; \codigo{|>} está siempre disponible.
\end{itemize}

En este libro usamos \codigo{\%>\%} porque es el que encontrarás en la
mayoría de los ejemplos en internet, pero puedes usar el que prefieras. Si
quieres que el atajo \tecla{Ctrl} + \tecla{Shift} + \tecla{M} escriba
\codigo{|>}, activa la casilla \emph{Use native pipe operator} en
\menu{Tools > Global Options > Code}.

\begin{hazlotu}
Reescribe con pipe la siguiente instrucción anidada y léela en voz alta
usando ``y luego'':
\begin{rnoejecutar}
head(arrange(filter(productos, categoria == "Ropa"), precio_catalogo), 3)
\end{rnoejecutar}
\end{hazlotu}

% ----------------------------------------------------------------------------
\section{Los verbos de \texorpdfstring{\paquete{dplyr}}{dplyr}}
\label{sec:m2-verbos}

\paquete{dplyr} está construido alrededor de unos cuantos \termino{verbos}:
funciones cuyo nombre dice lo que hacen con la tabla. Todos siguen las mismas
reglas, y una vez que las entiendes, aprender un verbo nuevo es trivial:

\begin{enumerate}
  \item El primer argumento siempre es un \emph{data frame} (por eso encajan
        tan bien con el pipe).
  \item Los demás argumentos usan los nombres de las columnas \emph{sin
        comillas} y sin \codigo{\$}: escribes \codigo{total}, no
        \codigo{ventas\$total}.
  \item El resultado siempre es un \emph{data frame} nuevo. La tabla original
        \emph{no se modifica}; si quieres conservar el resultado, guárdalo con
        \codigo{<-}.
\end{enumerate}

\begin{table}[H]
\centering
\caption{Los seis verbos fundamentales de \paquete{dplyr}.}
\label{tab:m2-verbos}
\small
\begin{tabularx}{\textwidth}{@{}lXl@{}}
\toprule
Verbo & Qué hace & Actúa sobre \\
\midrule
\codigo{filter()}    & Se queda con las filas que cumplen una condición & Filas \\
\codigo{select()}    & Elige, quita o reordena columnas & Columnas \\
\codigo{mutate()}    & Crea columnas nuevas o modifica las existentes & Columnas \\
\codigo{arrange()}   & Ordena las filas & Filas \\
\codigo{summarise()} & Resume muchas filas en una (totales, promedios\ldots) & Filas \\
\codigo{group\_by()} & Divide la tabla en grupos para que los demás verbos
                       trabajen grupo por grupo & Grupos \\
\bottomrule
\end{tabularx}
\end{table}

% ............................................................................
\subsection{\texorpdfstring{\codigo{filter()}}{filter()}: quedarte con las filas que importan}
\label{sec:m2-filter}

\codigo{filter()} recibe una o más condiciones lógicas y conserva solo las
filas donde la condición es \codigo{TRUE}. Las condiciones se construyen con
los operadores de la tabla~\ref{tab:m2-operadores}.

\begin{table}[H]
\centering
\caption{Operadores para construir condiciones en \codigo{filter()}.}
\label{tab:m2-operadores}
\small
\begin{tabularx}{\textwidth}{@{}lXl@{}}
\toprule
Operador & Significado & Ejemplo \\
\midrule
\codigo{==}  & Igual a (¡dos signos!) & \codigo{tienda\_id == 3} \\
\codigo{!=}  & Distinto de & \codigo{trimestre != "Q4"} \\
\codigo{>}, \codigo{>=} & Mayor, mayor o igual & \codigo{total >= 1000} \\
\codigo{<}, \codigo{<=} & Menor, menor o igual & \codigo{cantidad < 3} \\
\codigo{\&} o coma & Y: se deben cumplir ambas & \codigo{tienda\_id == 1, cantidad > 5} \\
\codigo{|}   & O: basta con que se cumpla una & \codigo{descuento\_pct == 20 | cantidad == 10} \\
\codigo{!}   & No (niega la condición) & \codigo{!activo} \\
\codigo{\%in\%} & Pertenece a una lista de valores & \codigo{tienda\_id \%in\% c(1, 4)} \\
\codigo{between()} & Dentro de un rango (incluye extremos) & \codigo{between(total, 500, 900)} \\
\codigo{is.na()} & Es un valor faltante & \codigo{is.na(email)} \\
\bottomrule
\end{tabularx}
\end{table}

\begin{ejemplo}{Ventas grandes}
El gerente de riesgos quiere revisar las ventas de más de \$4,500, porque
un ticket tan alto es poco común. ¿Cuántas son y de qué tiendas?
\end{ejemplo}

\begin{rcode}
ventas %>%
  filter(total > 4500) %>%                       # condición numérica
  select(fecha, tienda_id, cantidad, precio_unitario, total)
\end{rcode}
\begin{rsalida}
# A tibble: 4 x 5
  fecha      tienda_id cantidad precio_unitario total
  <date>         <dbl>    <dbl>           <dbl> <dbl>
1 2023-02-15         8       10            488. 4638.
2 2023-04-21         2       10            492. 4925.
3 2023-07-24         8       10            495. 4950.
4 2023-09-04         6       10            487. 4875.
\end{rsalida}

Solo 4 ventas del año superan los \$4,500. Fíjate en que las cuatro son de
10 unidades (el máximo posible) con precios unitarios cercanos a \$490: no
parecen errores sino compras al mayoreo de productos caros.

\begin{ejemplo}{Varias condiciones a la vez (Y)}
¿Qué ventas hizo la tienda 9 (la sucursal \emph{Premium}) en el primer
trimestre con algún descuento?
\end{ejemplo}

Separar condiciones con comas equivale a unirlas con \codigo{\&}: se deben
cumplir \emph{todas}.

\begin{rcode}
ventas %>%
  filter(tienda_id == 9,            # tienda Premium
         trimestre == "Q1",         # y primer trimestre
         descuento_pct > 0) %>%     # y con descuento
  select(fecha, producto_id, descuento_pct, total)
\end{rcode}
\begin{rsalida}
# A tibble: 5 x 4
  fecha      producto_id descuento_pct total
  <date>           <dbl>         <dbl> <dbl>
1 2023-01-29          39            10 3650.
2 2023-02-14          25             5 1468.
3 2023-02-27          15             5  155.
4 2023-03-16          49             5 1290.
5 2023-03-18           4            10 1486.
\end{rsalida}

\begin{ejemplo}{Una u otra condición (O) y listas de valores}
Mercadotecnia quiere las ventas ``extremas'': las que tuvieron el descuento
máximo (20\%) \emph{o} se llevaron 10 unidades. Y, aparte, las ventas de
las cuatro tiendas de la CDMX (tiendas 1, 4, 6 y 9).
\end{ejemplo}

\begin{rcode}
# O: basta con que se cumpla una de las dos condiciones
ventas %>%
  filter(descuento_pct == 20 | cantidad == 10) %>%
  nrow()

# %in%: la tienda está en la lista (más corto que cuatro "|")
ventas_cdmx <- ventas %>%
  filter(tienda_id %in% c(1, 4, 6, 9))
nrow(ventas_cdmx)
\end{rcode}
\begin{rsalida}
[1] 54
[1] 135
\end{rsalida}

54 ventas cumplen al menos una de las dos condiciones ``extremas'', y
135 ventas (más de un tercio del total) se hicieron en tiendas de la
CDMX. Usa \codigo{\%in\%} siempre que compares contra una lista: es más
legible y evitas errores como escribir
\codigo{tienda\_id == 1 | 4}, que \emph{no} significa lo que parece.

\begin{ejemplo}{Rangos de montos y de fechas}
¿Cuántas ventas están en el rango medio, entre \$1,000 y \$2,000? ¿Qué se
vendió durante El Buen Fin 2023 (del 17 al 20 de noviembre)?
\end{ejemplo}

\begin{rcode}
# between() incluye ambos extremos: 1000 <= total <= 2000
ventas %>%
  filter(between(total, 1000, 2000)) %>%
  nrow()

# Con fechas: compara contra objetos Date, no contra texto
ventas %>%
  filter(fecha >= as.Date("2023-11-17"),
         fecha <= as.Date("2023-11-20")) %>%
  select(fecha, dia_semana, producto_id, cantidad, total)
\end{rcode}
\begin{rsalida}
[1] 96
# A tibble: 4 x 5
  fecha      dia_semana producto_id cantidad total
  <date>     <chr>            <dbl>    <dbl> <dbl>
1 2023-11-17 viernes             39        7 1731.
2 2023-11-18 sábado              30        7 2939.
3 2023-11-19 domingo             31        5  142.
4 2023-11-20 lunes               30        8 3358.
\end{rsalida}

96 ventas caen entre \$1,000 y \$2,000. Durante El Buen Fin hubo 4 ventas
(recuerda que este archivo tiene una venta por día). Como la columna
\codigo{fecha} es de tipo \texttt{<date>}, la comparamos con
\codigo{as.Date("2023-11-17")}; así R compara fechas reales y no texto.

\begin{ejemplo}{Negaciones y valores faltantes}
¿Qué productos del catálogo están inactivos? Y en una tabla de metas de
vendedores donde algunas metas aún no se capturan, ¿quiénes no tienen meta?
\end{ejemplo}

La columna \codigo{activo} ya es lógica (\codigo{TRUE}/\codigo{FALSE}), así
que basta con negarla con \codigo{!}:

\begin{rcode}
productos %>%
  filter(!activo) %>%                 # los que NO están activos
  select(producto_id, nombre_producto, categoria, stock_actual)
\end{rcode}
\begin{rsalida}
# A tibble: 7 x 4
  producto_id nombre_producto categoria  stock_actual
        <dbl> <chr>           <chr>             <dbl>
1           7 Producto 07     Deportes            123
2          13 Producto 13     Ropa                177
3          15 Producto 15     Mascotas            455
4          25 Producto 25     Juguetes            429
5          38 Producto 38     Alimentos           141
6          39 Producto 39     Alimentos             8
7          46 Producto 46     Automotriz          118
\end{rsalida}

Hay 7 productos inactivos, y casi todos tienen todavía más de 100
unidades en almacén (el \emph{Producto 15}, 455): inventario inmovilizado,
un buen tema para la junta de compras.

Para ilustrar los faltantes creamos una tabla pequeña con
\codigo{tibble()}. Cada argumento es una columna:

\begin{rcode}
metas <- tibble(
  vendedor = c("Ana", "Luis", "Rosa", "Jorge", "Elena"),
  meta     = c(50000, NA, 42000, NA, 61000)   # NA = aún no capturada
)
# filter() descarta en silencio las filas donde la condición es NA
metas %>% filter(meta > 45000)
# Para encontrar los faltantes usa is.na(), nunca meta == NA
metas %>% filter(is.na(meta))
\end{rcode}
\begin{rsalida}
# A tibble: 2 x 2
  vendedor  meta
  <chr>    <dbl>
1 Ana      50000
2 Elena    61000
# A tibble: 2 x 2
  vendedor  meta
  <chr>    <dbl>
1 Luis        NA
2 Jorge       NA
\end{rsalida}

\begin{advertencia}
\codigo{meta == NA} no funciona: cualquier comparación con \codigo{NA} da
\codigo{NA} (``no sé''), y \codigo{filter()} descarta esas filas. Por eso
Luis y Jorge no aparecieron en el primer filtro \emph{ni aparecerían} con
\codigo{meta <= 45000}. Si necesitas conservarlos, sé explícito:
\codigo{filter(meta > 45000 | is.na(meta))}.
\end{advertencia}

\begin{ejemplo}{Filtrar por un patrón de texto}
En \codigo{transacciones}, ¿cuántas se pagaron con algún tipo de tarjeta?
\end{ejemplo}

La función \codigo{str\_detect()} de \paquete{stringr} devuelve
\codigo{TRUE} si el texto contiene el patrón buscado. La veremos a fondo en
la sección~\ref{sec:m2-stringr}.

\begin{rcode}
transacciones %>%
  filter(str_detect(metodo_pago, "Tarjeta")) %>%  # contiene "Tarjeta"
  count(metodo_pago)                              # cuántas de cada una
\end{rcode}
\begin{rsalida}
# A tibble: 2 x 2
  metodo_pago         n
  <chr>           <int>
1 Tarjeta Crédito   350
2 Tarjeta Débito    233
\end{rsalida}

583 de las 1000 transacciones (350 + 233) se pagaron con tarjeta.
La función \codigo{count()} cuenta filas por categoría; la estudiaremos en
la sección~\ref{sec:m2-summarise}.

\begin{hazlotu}
Con \codigo{productos}: (a) filtra los de la categoría ``Electrónica'' con
precio de catálogo mayor a \$400; (b) filtra los que están en
``Stock Bajo'' o ``Sin Stock'' usando \codigo{\%in\%}; (c) cuenta cuántos
clientes de \codigo{clientes} son premium y viven en Monterrey.
\end{hazlotu}

% ............................................................................
\subsection{\texorpdfstring{\codigo{select()}}{select()}: elegir, quitar y renombrar columnas}
\label{sec:m2-select}

Las tablas reales tienen decenas de columnas y casi nunca las necesitas
todas. \codigo{select()} se queda con las que le pidas, en el orden en que
las pidas.

\begin{ejemplo}{Una lista de precios para el equipo de ventas}
El equipo de piso solo necesita el nombre, la categoría y el precio de cada
producto; el costo es confidencial y no debe aparecer.
\end{ejemplo}

\begin{rcode}
lista_precios <- productos %>%
  select(nombre_producto, categoria, precio_catalogo)   # por nombre
head(lista_precios, 4)
\end{rcode}
\begin{rsalida}
# A tibble: 4 x 3
  nombre_producto categoria  precio_catalogo
  <chr>           <chr>                <dbl>
1 Producto 01     Juguetes              300.
2 Producto 02     Belleza               313.
3 Producto 03     Automotriz            400.
4 Producto 04     Juguetes              744.
\end{rsalida}

También puedes seleccionar un \emph{rango} de columnas consecutivas con
\codigo{:}, o \emph{quitar} columnas con un signo menos:

\begin{rcode}
# Rango: de fecha hasta cantidad (columnas consecutivas)
ventas %>% select(fecha:cantidad) %>% names()

# Quitar: todo menos las etiquetas de texto que no usaremos
ventas %>% select(-dia_semana, -mes, -trimestre) %>% names()
\end{rcode}
\begin{rsalida}
[1] "fecha"       "dia_semana"  "producto_id" "cliente_id" 
[5] "cantidad"   
 [1] "fecha"           "producto_id"     "cliente_id"     
 [4] "cantidad"        "precio_unitario" "descuento_pct"  
 [7] "tienda_id"       "vendedor_id"     "subtotal"       
[10] "descuento_monto" "total"          
\end{rsalida}

La función \codigo{names()} devuelve los nombres de las columnas; aquí la
usamos para mostrar el resultado de forma compacta.

\subsubsection{Funciones auxiliares de selección}

Cuando las columnas siguen una convención de nombres (como
\texttt{\_id} al final de las llaves), las \termino{funciones auxiliares}
(\emph{helpers}) te ahorran escribirlas una por una:

\begin{rcode}
# Todas las llaves: terminan en "_id"
ventas %>% select(ends_with("_id")) %>% names()

# Columnas de inventario: empiezan con "stock"
productos %>% select(producto_id, starts_with("stock")) %>% head(3)

# Columnas que contienen "fecha" en cualquier parte del nombre
clientes %>% select(cliente_id, contains("fecha")) %>% head(3)

# Solo las columnas numéricas (útil antes de calcular resúmenes)
productos %>% select(where(is.numeric)) %>% names()
\end{rcode}
\begin{rsalida}
[1] "producto_id" "cliente_id"  "tienda_id"   "vendedor_id"
# A tibble: 3 x 3
  producto_id stock_actual stock_minimo
        <dbl>        <dbl>        <dbl>
1           1          374           21
2           2          371           46
3           3           22           26
# A tibble: 3 x 2
  cliente_id fecha_registro
       <dbl> <date>        
1          1 2023-10-12    
2          2 2020-01-30    
3          3 2020-10-17    
[1] "producto_id"     "precio_catalogo" "costo"          
[4] "stock_actual"    "stock_minimo"    "margen_pct"     
\end{rsalida}

\codigo{where()} recibe una función que devuelve \codigo{TRUE} o
\codigo{FALSE} para cada columna: \codigo{is.numeric}, \codigo{is.character},
\codigo{is.logical}\ldots{} Otras auxiliares útiles son
\codigo{everything()} (``todas las demás''), \codigo{matches()} (patrón de
texto) y \codigo{any\_of()} (una lista de nombres que pueden no existir).

\subsubsection{Renombrar y reacomodar: \texorpdfstring{\codigo{rename()}}{rename()} y \texorpdfstring{\codigo{relocate()}}{relocate()}}

\codigo{rename()} cambia nombres sin quitar columnas; la sintaxis es
\codigo{nuevo = viejo}. \codigo{relocate()} mueve columnas de lugar sin
quitar ninguna.

\begin{rcode}
tiendas %>%
  rename(tienda = nombre_tienda, m2 = tamano_m2) %>%  # nuevo = viejo
  select(tienda_id, tienda, ciudad, m2) %>%
  head(3)

# Poner el total justo después de la fecha para leerlo más fácil
ventas %>%
  relocate(total, .after = fecha) %>%
  head(3)
\end{rcode}
\begin{rsalida}
# A tibble: 3 x 4
  tienda_id tienda          ciudad         m2
      <dbl> <chr>           <chr>       <dbl>
1         1 Sucursal Centro CDMX          500
2         2 Sucursal Norte  Guadalajara   450
3         3 Sucursal Sur    Monterrey     600
# A tibble: 3 x 14
  fecha      total dia_semana producto_id cliente_id cantidad
  <date>     <dbl> <chr>            <dbl>      <dbl>    <dbl>
1 2023-01-01 3333. domingo             31         94       10
2 2023-01-02 2262  lunes               15        169        8
3 2023-01-03  801. martes              14        136        3
# i 8 more variables: precio_unitario <dbl>, descuento_pct <dbl>,
#   tienda_id <dbl>, vendedor_id <dbl>, subtotal <dbl>,
#   descuento_monto <dbl>, mes <chr>, trimestre <chr>
\end{rsalida}

\begin{consejo}
Puedes renombrar dentro de \codigo{select()}:
\codigo{select(tienda = nombre\_tienda, ciudad)} elige \emph{y} renombra en
un paso. Y para mandar una columna al principio sin perder las demás, usa
\codigo{select(total, everything())} o \codigo{relocate(total)}.
\end{consejo}

\begin{hazlotu}
De \codigo{empleados}, arma una tabla con \codigo{empleado\_id}, el nombre
renombrado como \codigo{vendedor}, y todas las columnas que empiecen con
``salario''. Después usa \codigo{select(where(is.character))} para ver qué
columnas de texto tiene la tabla.
\end{hazlotu}

% ............................................................................
\subsection{\texorpdfstring{\codigo{mutate()}}{mutate()}: crear columnas calculadas}
\label{sec:m2-mutate}

\codigo{mutate()} agrega columnas nuevas (o reemplaza existentes) calculadas
a partir de las demás. Es el equivalente a escribir una fórmula en una
columna de Excel y arrastrarla hacia abajo, solo que se aplica a todas las
filas de una vez y queda documentada en el código.

\begin{ejemplo}{IVA y total con impuestos}
Los montos de \codigo{ventas} no incluyen IVA. Contabilidad necesita, para
cada venta, el IVA (16\%) y el total facturado.
\end{ejemplo}

\begin{rcode}
ventas %>%
  mutate(
    iva            = total * 0.16,        # 16% sobre el total
    total_con_iva  = total + iva          # usa la columna recién creada
  ) %>%
  select(fecha, subtotal, descuento_monto, total, iva, total_con_iva) %>%
  head(4)
\end{rcode}
\begin{rsalida}
# A tibble: 4 x 6
  fecha      subtotal descuento_monto total   iva total_con_iva
  <date>        <dbl>           <dbl> <dbl> <dbl>         <dbl>
1 2023-01-01    3508.           175.  3333. 533.          3866.
2 2023-01-02    2262              0   2262  362.          2624.
3 2023-01-03     890.            89.0  801. 128.           929.
4 2023-01-04     445.             0    445.  71.2          516.
\end{rsalida}

Observa dos cosas. Dentro del mismo \codigo{mutate()} puedes usar una
columna que acabas de crear (\codigo{total\_con\_iva} usa \codigo{iva}). Y la
primera venta: subtotal de \$3,507.90 menos \$175.40 de descuento da
\$3,332.50; con IVA, \$3,865.71.

\begin{ejemplo}{¿Qué tan rentable es cada producto?}
Con el catálogo, calcula el margen bruto en pesos y en porcentaje, y
ordena del menos al más rentable.
\end{ejemplo}

\begin{rcode}
productos %>%
  mutate(
    margen_pesos = precio_catalogo - costo,
    margen_calc  = round(margen_pesos / precio_catalogo * 100, 1)
  ) %>%
  select(nombre_producto, precio_catalogo, costo,
         margen_pesos, margen_calc) %>%
  arrange(margen_calc) %>%            # del menor al mayor margen
  head(5)
\end{rcode}
\begin{rsalida}
# A tibble: 5 x 5
  nombre_producto precio_catalogo costo margen_pesos margen_calc
  <chr>                     <dbl> <dbl>        <dbl>       <dbl>
1 Producto 05                25.5  355.       -329.      -1292. 
2 Producto 46                64.5  178.       -113.       -175. 
3 Producto 32               143.   210.        -67.0       -46.7
4 Producto 26               115.   159.        -44.2       -38.3
5 Producto 29               298.   392.        -93.9       -31.5
\end{rsalida}

\begin{negocio}{Tu primer hallazgo de calidad de datos}
¡Hay productos que cuestan más de lo que se venden! El \emph{Producto 05}
tiene un precio de catálogo de \$25.47 y un costo de \$354.63: un margen de
$-1292$\%. Puede ser un error de captura (¿el precio debía ser \$254.70?) o un
producto que se vende en promoción. Como analista, tu trabajo no es ocultarlo
sino \emph{reportarlo}: en la sección~\ref{sec:m2-caso} contaremos cuántos
productos están en esta situación.
\end{negocio}

\subsubsection{Condiciones: \texorpdfstring{\codigo{if\_else()}}{if\_else()} y \texorpdfstring{\codigo{case\_when()}}{case\_when()}}

\codigo{if\_else(condición, valor\_si\_verdadero, valor\_si\_falso)} crea
una columna con dos posibles valores, como la función \codigo{SI()} de
Excel:

\begin{rcode}
ventas %>%
  mutate(tipo_dia = if_else(dia_semana %in% c("sábado", "domingo"),
                            "Fin de semana", "Entre semana")) %>%
  count(tipo_dia)
\end{rcode}
\begin{rsalida}
# A tibble: 2 x 2
  tipo_dia          n
  <chr>         <int>
1 Entre semana    260
2 Fin de semana   105
\end{rsalida}

Cuando hay más de dos opciones, anidar \codigo{if\_else()} se vuelve
ilegible. Para eso existe \codigo{case\_when()}: una lista de reglas
\codigo{condición \textasciitilde{} valor} que se evalúan \emph{en orden}; a cada
fila se le asigna el valor de la \emph{primera} regla que cumple.

\begin{ejemplo}{Segmentar ventas por tamaño de ticket}
Mercadotecnia quiere clasificar cada venta en cuatro segmentos: Chico
(menos de \$500), Mediano (\$500 a \$1,499), Grande (\$1,500 a \$2,999) y
Premium (\$3,000 o más). ¿Cuántas ventas y cuánto dinero hay en cada uno?
\end{ejemplo}

\begin{rcode}
ventas_segmentadas <- ventas %>%
  mutate(segmento_ticket = case_when(
    total < 500   ~ "1. Chico",       # se evalúa primero
    total < 1500  ~ "2. Mediano",     # solo si no fue Chico
    total < 3000  ~ "3. Grande",
    .default      = "4. Premium"      # todo lo demás
  ))

ventas_segmentadas %>%
  group_by(segmento_ticket) %>%
  summarise(ventas = n(), ingresos = sum(total))
\end{rcode}
\begin{rsalida}
# A tibble: 4 x 3
  segmento_ticket ventas ingresos
  <chr>            <int>    <dbl>
1 1. Chico            98   25588.
2 2. Mediano         129  125261.
3 3. Grande          101  227370.
4 4. Premium          37  135697.
\end{rsalida}

Una lectura de negocio: los tickets Premium son solo 37 ventas (una de
cada diez), pero generan \$135,697, más de cinco veces lo que dejan los 98
tickets Chicos juntos (\$25,588). El segmento que más dinero aporta es el
Grande (\$227,370). La estrategia comercial debe cuidar a los tickets
altos. (Adelantamos
\codigo{group\_by()} y \codigo{summarise()}, que veremos en detalle muy
pronto.)

\begin{consejo}
Pon un número al inicio de cada etiqueta (``1. Chico'', ``2. Mediano''\ldots)
y los grupos aparecerán en orden lógico y no alfabético. El argumento
\codigo{.default} (dplyr 1.1 o posterior) asigna un valor a todo lo que no
cumplió ninguna regla; en código antiguo verás \codigo{TRUE \textasciitilde{} "Premium"},
que hace lo mismo.
\end{consejo}

\subsubsection{Varias columnas a la vez: \texorpdfstring{\codigo{across()}}{across()}}

Si quieres aplicar la misma operación a varias columnas, no repitas el
código: usa \codigo{across(columnas, función)}. Las columnas se eligen con
la misma sintaxis de \codigo{select()}.

\begin{rcode}
# Redondear a 0 decimales las tres columnas de dinero
ventas %>%
  mutate(across(c(subtotal, descuento_monto, total), round)) %>%
  select(fecha, subtotal, descuento_monto, total) %>%
  head(3)

# Con una función "al vuelo": ~ .x indica "cada columna"
productos %>%
  mutate(across(c(precio_catalogo, costo), ~ .x * 1.16,
                .names = "{.col}_con_iva")) %>%
  select(nombre_producto, ends_with("_con_iva")) %>%
  head(3)
\end{rcode}
\begin{rsalida}
# A tibble: 3 x 4
  fecha      subtotal descuento_monto total
  <date>        <dbl>           <dbl> <dbl>
1 2023-01-01     3508             175  3333
2 2023-01-02     2262               0  2262
3 2023-01-03      890              89   801
# A tibble: 3 x 3
  nombre_producto precio_catalogo_con_iva costo_con_iva
  <chr>                             <dbl>         <dbl>
1 Producto 01                        348.          338.
2 Producto 02                        363.          463.
3 Producto 03                        464.          341.
\end{rsalida}

La notación \codigo{\textasciitilde{} .x * 1.16} es una \termino{función anónima}: una
fórmula corta donde \codigo{.x} representa cada columna. El argumento
\codigo{.names} controla el nombre de las columnas nuevas
(\codigo{\{.col\}} es el nombre original); sin él, \codigo{across()}
sobrescribe las columnas originales.

\subsubsection{Conservar solo lo calculado: \texorpdfstring{\codigo{.keep}}{.keep} y \texorpdfstring{\codigo{transmute()}}{transmute()}}

Por defecto, \codigo{mutate()} conserva todas las columnas. Con el argumento
\codigo{.keep} decides qué más se queda: \codigo{"used"} (solo las columnas
usadas en el cálculo) o \codigo{"none"} (solo las nuevas). La antigua función
\codigo{transmute()} equivale a \codigo{.keep = "none"}.

\begin{rcode}
# Solo las columnas que participaron en el cálculo, más la nueva
ventas %>%
  mutate(precio_neto = total / cantidad, .keep = "used") %>%
  head(3)

# Solo las nuevas (equivalente a transmute())
ventas %>%
  mutate(anio = str_sub(mes, 1, 4), ticket = round(total),
         .keep = "none") %>%
  head(3)
\end{rcode}
\begin{rsalida}
# A tibble: 3 x 3
  cantidad total precio_neto
     <dbl> <dbl>       <dbl>
1       10 3333.        333.
2        8 2262         283.
3        3  801.        267.
# A tibble: 3 x 2
  anio  ticket
  <chr>  <dbl>
1 2023    3333
2 2023    2262
3 2023     801
\end{rsalida}

\begin{advertencia}
\codigo{mutate()} no modifica \codigo{ventas}: devuelve una copia con la
columna nueva. Si ejecutas \codigo{ventas \%>\% mutate(iva = total * 0.16)}
y luego buscas \codigo{ventas\$iva}, no existirá. Para conservarla escribe
\codigo{ventas <- ventas \%>\% mutate(\ldots)} o, mejor, guarda el resultado
en un objeto con otro nombre y deja intacta la tabla original.
\end{advertencia}

\begin{hazlotu}
En \codigo{clientes}, crea una columna \codigo{rango\_edad} con
\codigo{case\_when()}: ``18--29'', ``30--44'', ``45--59'' y ``60+''. Cuenta
cuántos clientes hay en cada rango. Pista: termina con
\codigo{count(rango\_edad)}.
\end{hazlotu}

% ............................................................................
\subsection{\texorpdfstring{\codigo{arrange()}}{arrange()} y la familia \texorpdfstring{\codigo{slice\_*()}}{slice\_*()}: ordenar y quedarse con el top}
\label{sec:m2-arrange}

\codigo{arrange()} ordena las filas de menor a mayor (orden ascendente). Para
ordenar de mayor a menor, envuelve la columna en \codigo{desc()}. Si le das
varias columnas, ordena por la primera y usa las siguientes para desempatar.

\begin{ejemplo}{Los productos más caros y las tiendas por ciudad}
¿Cuáles son los cinco productos con mayor precio de catálogo? ¿Cómo quedan
las tiendas ordenadas por ciudad y, dentro de cada ciudad, de la más grande
a la más chica?
\end{ejemplo}

\begin{rcode}
productos %>%
  arrange(desc(precio_catalogo)) %>%    # de mayor a menor
  select(nombre_producto, categoria, precio_catalogo) %>%
  head(5)

tiendas %>%
  arrange(ciudad, desc(tamano_m2)) %>%  # ciudad A-Z; luego m2 desc.
  select(ciudad, nombre_tienda, tamano_m2)
\end{rcode}
\begin{rsalida}
# A tibble: 5 x 3
  nombre_producto categoria   precio_catalogo
  <chr>           <chr>                 <dbl>
1 Producto 28     Mascotas               774.
2 Producto 47     Automotriz             766.
3 Producto 41     Electrónica            757.
4 Producto 04     Juguetes               744.
5 Producto 20     Belleza                742.
# A tibble: 10 x 3
   ciudad      nombre_tienda    tamano_m2
   <chr>       <chr>                <dbl>
 1 CDMX        Sucursal Premium       800
 2 CDMX        Sucursal Plaza A       700
 3 CDMX        Sucursal Centro        500
 4 CDMX        Sucursal Este          400
 5 Guadalajara Sucursal Plaza B       480
 6 Guadalajara Sucursal Norte         450
 7 Monterrey   Sucursal Sur           600
 8 Puebla      Sucursal Oeste         550
 9 Querétaro   Sucursal Express       300
10 Tijuana     Sucursal Outlet        350
\end{rsalida}

La CDMX tiene cuatro tiendas y la más grande es la \emph{Sucursal Premium},
con 800~m$^2$. Observa que el texto se ordena alfabéticamente respetando los
acentos del español (según la configuración regional de tu computadora).

Para quedarte con las primeras o las últimas filas existe la familia
\codigo{slice\_*()}:

\begin{itemize}
  \item \codigo{slice\_head(n = 3)} y \codigo{slice\_tail(n = 3)}: las
        primeras o las últimas 3 filas, en el orden actual.
  \item \codigo{slice\_max(columna, n = 5)} y
        \codigo{slice\_min(columna, n = 5)}: las 5 filas con los valores más
        altos o más bajos de una columna. No necesitas ordenar antes.
  \item \codigo{slice\_sample(n = 10)}: 10 filas al azar (útil para
        auditorías; recuerda fijar la semilla con \codigo{set.seed()}).
\end{itemize}

\begin{rcode}
# Top 3 ventas del año por monto
ventas %>%
  slice_max(total, n = 3) %>%
  select(fecha, tienda_id, total)

# Los 3 productos más baratos
productos %>%
  slice_min(precio_catalogo, n = 3) %>%
  select(nombre_producto, precio_catalogo)

# Una muestra aleatoria de 3 ventas para auditar
set.seed(123)
ventas %>%
  slice_sample(n = 3) %>%
  select(fecha, producto_id, total)
\end{rcode}
\begin{rsalida}
# A tibble: 3 x 3
  fecha      tienda_id total
  <date>         <dbl> <dbl>
1 2023-07-24         8 4950.
2 2023-04-21         2 4925.
3 2023-09-04         6 4875.
# A tibble: 3 x 2
  nombre_producto precio_catalogo
  <chr>                     <dbl>
1 Producto 05                25.5
2 Producto 38                36.6
3 Producto 40                50.1
# A tibble: 3 x 3
  fecha      producto_id total
  <date>           <dbl> <dbl>
1 2023-06-28          20 1993.
2 2023-01-14          27  982.
3 2023-07-14          14  897.
\end{rsalida}

\subsubsection{Top-N por grupo}

La verdadera potencia de \codigo{slice\_max()} aparece cuando la combinas con
\codigo{group\_by()}: obtienes el top-N \emph{de cada grupo} en una sola
instrucción. En Excel esto requiere fórmulas complicadas o filtros manuales
grupo por grupo.

\begin{ejemplo}{La mejor venta de cada trimestre y el producto más caro de cada categoría}
El director quiere, para su presentación, la venta más grande de cada
trimestre. Compras quiere saber cuál es el producto más caro de cada
categoría.
\end{ejemplo}

\begin{rcode}
# La venta más grande de cada trimestre
ventas %>%
  group_by(trimestre) %>%
  slice_max(total, n = 1) %>%
  ungroup() %>%
  select(trimestre, fecha, tienda_id, total)

# El producto más caro de cada categoría
productos %>%
  group_by(categoria) %>%
  slice_max(precio_catalogo, n = 1) %>%
  ungroup() %>%
  select(categoria, nombre_producto, precio_catalogo)
\end{rcode}
\begin{rsalida}
# A tibble: 4 x 4
  trimestre fecha      tienda_id total
  <chr>     <date>         <dbl> <dbl>
1 Q1        2023-02-15         8 4638.
2 Q2        2023-04-21         2 4925.
3 Q3        2023-07-24         8 4950.
4 Q4        2023-12-01         7 4255.
# A tibble: 10 x 3
   categoria   nombre_producto precio_catalogo
   <chr>       <chr>                     <dbl>
 1 Alimentos   Producto 33                683.
 2 Automotriz  Producto 47                766.
 3 Belleza     Producto 20                742.
 4 Deportes    Producto 19                583.
 5 Electrónica Producto 41                757.
 6 Hogar       Producto 10                691.
 7 Juguetes    Producto 04                744.
 8 Libros      Producto 35                696.
 9 Mascotas    Producto 28                774.
10 Ropa        Producto 42                368.
\end{rsalida}

\codigo{group\_by()} parte la tabla en grupos y \codigo{slice\_max()} se
aplica dentro de cada uno. Al final, \codigo{ungroup()} quita la agrupación
para que no afecte los pasos siguientes (lo explicamos en la próxima
sección). Si hubiera empates, \codigo{slice\_max()} devolvería todas las
filas empatadas; agrega \codigo{with\_ties = FALSE} si necesitas exactamente
\codigo{n} filas.

\begin{hazlotu}
Obtén las 2 ventas con más unidades (\codigo{cantidad}) de \emph{cada}
tienda. ¿Cuántas filas esperas? ¿Cuántas obtienes? Si no coinciden,
explica por qué y prueba con \codigo{with\_ties = FALSE}.
\end{hazlotu}

% ............................................................................
\subsection{\texorpdfstring{\codigo{summarise()}}{summarise()} y \texorpdfstring{\codigo{group\_by()}}{group\_by()}: de miles de filas a un KPI}
\label{sec:m2-summarise}

\codigo{summarise()} (o \codigo{summarize()}, ambas funcionan) reduce muchas
filas a una sola fila de resultados. Dentro usas \termino{funciones de
resumen}: reciben un vector de muchos valores y devuelven uno solo.

\begin{table}[H]
\centering
\caption{Funciones de resumen más usadas en BI.}
\label{tab:m2-resumen-fun}
\small
\begin{tabularx}{\textwidth}{@{}lX@{}}
\toprule
Función & Qué calcula \\
\midrule
\codigo{n()} & Número de filas (transacciones, registros) \\
\codigo{n\_distinct(x)} & Número de valores distintos (clientes únicos, productos vendidos) \\
\codigo{sum(x)} & Suma (ingresos, unidades) \\
\codigo{mean(x)}, \codigo{median(x)} & Promedio y mediana (ticket promedio, ticket típico) \\
\codigo{min(x)}, \codigo{max(x)} & Mínimo y máximo (primera y última compra, venta récord) \\
\codigo{sd(x)} & Desviación estándar (qué tan variables son los valores) \\
\codigo{first(x)}, \codigo{last(x)} & Primer y último valor según el orden actual \\
\bottomrule
\end{tabularx}
\end{table}

\begin{ejemplo}{El tablero de KPIs del año}
La dirección quiere los indicadores clave de 2023 en una sola tabla:
ingresos, número de ventas, unidades, ticket promedio y mediano, venta mínima
y máxima, y clientes únicos.
\end{ejemplo}

\begin{rcode}
kpis_2023 <- ventas %>%
  summarise(
    ingresos        = sum(total),
    num_ventas      = n(),                    # cuenta filas
    unidades        = sum(cantidad),
    ticket_promedio = mean(total),
    ticket_mediano  = median(total),
    venta_minima    = min(total),
    venta_maxima    = max(total),
    clientes_unicos = n_distinct(cliente_id)  # clientes distintos
  )
glimpse(kpis_2023)
\end{rcode}
\begin{rsalida}
Rows: 1
Columns: 8
$ ingresos        <dbl> 513916.6
$ num_ventas      <int> 365
$ unidades        <dbl> 2057
$ ticket_promedio <dbl> 1407.991
$ ticket_mediano  <dbl> 1103.67
$ venta_minima    <dbl> 20.07
$ venta_maxima    <dbl> 4950.3
$ clientes_unicos <int> 172
\end{rsalida}

Usamos \codigo{glimpse()} porque la tabla tiene una sola fila y muchas
columnas. En 2023 se vendieron \$513,916.60 en 365 ventas (2,057 unidades).
El ticket promedio (\$1,408) es mayor que el mediano (\$1,103.67): unas cuantas
ventas muy grandes ``jalan'' el promedio hacia arriba, algo muy común en
ventas. Por eso en BI conviene reportar ambos. Además, 172 clientes
distintos compraron al menos una vez.

\subsubsection{Resumir por grupos}

Sola, \codigo{summarise()} da un total general. Precedida de
\codigo{group\_by()}, da un resumen \emph{por cada grupo}. Esta combinación
sigue la estrategia \termino{dividir-aplicar-combinar}
(figura~\ref{fig:m2-dac}): R divide la tabla según los valores de la columna
de agrupación, aplica el resumen a cada pedazo y junta los resultados en una
tabla nueva con una fila por grupo.

\begin{figure}[H]
\centering
\begin{tikzpicture}[
  celda/.style={draw=white, minimum width=1.25cm, minimum height=0.42cm,
                font=\ttfamily\scriptsize, inner sep=1pt},
  tit/.style={font=\sffamily\scriptsize\bfseries, text=gristexto},
  flecha/.style={-{Stealth[length=2.2mm]}, thick, color=acento}]
  % Tabla original
  \node[tit] at (0.62,0.55) {ventas};
  \foreach \tri/\val/\col [count=\i] in {Q1/100/primario,
      Q2/250/secundario, Q1/300/primario, Q2/50/secundario,
      Q1/200/primario}{
    \node[celda, fill=\col!25] at (0,-\i*0.42) {\tri};
    \node[celda, fill=\col!25] at (1.25,-\i*0.42) {\val};
  }
  % Dividir
  \node[tit] at (4.2,0.55) {dividir (group\_by)};
  \foreach \val [count=\i] in {100,300,200}{
    \node[celda, fill=primario!25] at (3.6,-\i*0.42) {Q1};
    \node[celda, fill=primario!25] at (4.85,-\i*0.42) {\val};
  }
  \foreach \val [count=\i] in {250,50}{
    \node[celda, fill=secundario!25] at (3.6,-1.5-\i*0.42) {Q2};
    \node[celda, fill=secundario!25] at (4.85,-1.5-\i*0.42) {\val};
  }
  % Aplicar
  \node[tit] at (8.1,0.55) {aplicar (sum)};
  \node[celda, fill=primario!25, minimum width=2.1cm] at (8.1,-0.84)
       {sum = 600};
  \node[celda, fill=secundario!25, minimum width=2.1cm] at (8.1,-2.34)
       {sum = 300};
  % Combinar
  \node[tit] at (11.6,0.55) {combinar};
  \node[celda, fill=gray!15, font=\ttfamily\scriptsize\bfseries]
       at (11.0,-0.84) {trim.};
  \node[celda, fill=gray!15, font=\ttfamily\scriptsize\bfseries]
       at (12.25,-0.84) {ingresos};
  \node[celda, fill=primario!25] at (11.0,-1.26) {Q1};
  \node[celda, fill=primario!25] at (12.25,-1.26) {600};
  \node[celda, fill=secundario!25] at (11.0,-1.68) {Q2};
  \node[celda, fill=secundario!25] at (12.25,-1.68) {300};
  % Flechas
  \draw[flecha] (2.05,-1.26) -- (2.8,-1.26);
  \draw[flecha] (5.65,-0.84) -- (6.9,-0.84);
  \draw[flecha] (5.65,-2.34) -- (6.9,-2.34);
  \draw[flecha] (9.3,-1.26) -- (10.2,-1.26);
\end{tikzpicture}
\caption{Dividir-aplicar-combinar: \codigo{group\_by(trimestre) \%>\%
summarise(ingresos = sum(total))}.}
\label{fig:m2-dac}
\end{figure}

\begin{ejemplo}{Ventas por trimestre}
¿Cómo se comportaron las ventas en cada trimestre? ¿Cuál fue el mejor?
\end{ejemplo}

\begin{rcode}
ventas %>%
  group_by(trimestre) %>%                 # un grupo por trimestre
  summarise(
    ingresos        = sum(total),
    num_ventas      = n(),
    ticket_promedio = mean(total)
  )
\end{rcode}
\begin{rsalida}
# A tibble: 4 x 4
  trimestre ingresos num_ventas ticket_promedio
  <chr>        <dbl>      <int>           <dbl>
1 Q1         126274.         90           1403.
2 Q2         120717.         91           1327.
3 Q3         134427.         92           1461.
4 Q4         132498.         92           1440.
\end{rsalida}

El tercer trimestre (julio a septiembre) fue el mejor, con \$134,427, y el
segundo el más flojo (\$120,717). Como todos los trimestres tienen casi el
mismo número de ventas (entre 90 y 92), la diferencia viene del ticket
promedio: \$1,461 en el Q3 contra \$1,327 en el Q2.

\begin{ejemplo}{Ranking de tiendas}
¿Qué tienda vende más? Queremos ingresos, número de ventas, unidades y
ticket promedio por tienda, de mayor a menor.
\end{ejemplo}

\begin{rcode}
ventas_tienda <- ventas %>%
  group_by(tienda_id) %>%
  summarise(
    ingresos        = sum(total),
    num_ventas      = n(),
    unidades        = sum(cantidad),
    ticket_promedio = round(mean(total), 2)
  ) %>%
  arrange(desc(ingresos))                 # la mejor tienda primero
ventas_tienda
\end{rcode}
\begin{rsalida}
# A tibble: 10 x 5
   tienda_id ingresos num_ventas unidades ticket_promedio
       <dbl>    <dbl>      <int>    <dbl>           <dbl>
 1         8   70155.         47      246           1493.
 2         9   61369.         46      260           1334.
 3         7   55290.         37      209           1494.
 4         2   54846.         38      198           1443.
 5         3   54066.         37      222           1461.
 6         6   54050.         34      219           1590.
 7         5   53611.         37      229           1449.
 8         1   43899.         32      169           1372.
 9        10   42171.         34      195           1240.
10         4   24459.         23      110           1063.
\end{rsalida}

La tienda 8 lidera con \$70,155 y la tienda 4 es la última con \$24,459:
apenas un tercio. El problema de la tienda 4 es doble: es la que menos
ventas hizo (23) \emph{y} la de menor ticket promedio (\$1,063). En cambio,
la tienda 6 tiene el ticket más alto (\$1,590) pero pocas ventas (34): su
reto es atraer más tráfico, no vender más caro. Este tipo de lectura, que
separa volumen (cuántas ventas) y valor (de cuánto cada una), es la esencia
del análisis de BI.

\subsubsection{Agrupar por varias columnas y el argumento \texorpdfstring{\codigo{.groups}}{.groups}}

Puedes agrupar por más de una columna. Observa el mensaje que aparece:

\begin{rcode}
ventas_trim_tienda <- ventas %>%
  group_by(trimestre, tienda_id) %>%
  summarise(ingresos = sum(total))
\end{rcode}
\begin{rsalida}
`summarise()` has grouped output by 'trimestre'. You can override
using the `.groups` argument.
\end{rsalida}

No es un error. \paquete{dplyr} te avisa que el resultado \emph{sigue
agrupado} por \codigo{trimestre}: cada \codigo{summarise()} ``pela'' solo el
último nivel de agrupación. Esto importa porque cualquier cálculo posterior
se hará por trimestre. Para no llevarte sorpresas, di explícitamente qué
quieres con \codigo{.groups}:

\begin{rcode}
ventas_trim_tienda <- ventas %>%
  group_by(trimestre, tienda_id) %>%
  summarise(ingresos = sum(total), .groups = "drop")  # sin grupos
ventas_trim_tienda
\end{rcode}
\begin{rsalida}
# A tibble: 40 x 3
   trimestre tienda_id ingresos
   <chr>         <dbl>    <dbl>
 1 Q1                1   16711.
 2 Q1                2   10435.
 3 Q1                3   11878.
 4 Q1                4    5750.
 5 Q1                5   14250.
 6 Q1                6   13988.
...
\end{rsalida}

\codigo{.groups = "drop"} quita todos los grupos (lo más seguro y lo que
usaremos casi siempre); \codigo{.groups = "keep"} los conserva todos. El
resultado tiene 40 filas: 4 trimestres $\times$ 10 tiendas.

\begin{consejo}
Desde \paquete{dplyr} 1.1 existe una alternativa más corta: el argumento
\codigo{.by} agrupa \emph{solo} para esa operación y devuelve siempre una
tabla sin grupos. \codigo{summarise(ingresos = sum(total), .by =
c(trimestre, tienda\_id))} equivale al código anterior. Verás ambas formas
en código real; en este libro usamos \codigo{group\_by()} porque es la más
difundida.
\end{consejo}

\subsubsection{Contar rápido con \texorpdfstring{\codigo{count()}}{count()}}

Contar filas por grupo es tan común que tiene su propio atajo:
\codigo{count(x)} equivale a \codigo{group\_by(x) \%>\% summarise(n = n())}.
Acepta \codigo{sort = TRUE} para ordenar de mayor a menor y \codigo{wt}
(\emph{weight}, peso) para \emph{sumar} una columna en lugar de contar filas.

\begin{rcode}
# ¿Cuántas transacciones por método de pago?
transacciones %>% count(metodo_pago, sort = TRUE)

# Con wt: suma el total en vez de contar filas
transacciones %>%
  count(metodo_pago, wt = total_transaccion, sort = TRUE,
        name = "ingresos")

# Dos columnas: combinaciones de tienda y método de pago
transacciones %>% count(tienda_id, metodo_pago) %>% head(5)
\end{rcode}
\begin{rsalida}
# A tibble: 4 x 2
  metodo_pago         n
  <chr>           <int>
1 Tarjeta Crédito   350
2 Efectivo          315
3 Tarjeta Débito    233
4 Transferencia     102
# A tibble: 4 x 2
  metodo_pago     ingresos
  <chr>              <dbl>
1 Tarjeta Crédito  351171.
2 Efectivo         303735.
3 Tarjeta Débito   220855.
4 Transferencia     84232.
# A tibble: 5 x 3
  tienda_id metodo_pago         n
      <dbl> <chr>           <int>
1         1 Efectivo           26
2         1 Tarjeta Crédito    37
3         1 Tarjeta Débito     14
4         1 Transferencia      10
5         2 Efectivo           36
\end{rsalida}

La tarjeta de crédito es el método más frecuente (350 transacciones) y
también el que más dinero mueve (\$351,171). El argumento \codigo{name}
cambia el nombre de la columna de conteo, que por defecto se llama
\codigo{n}.

\subsubsection{Porcentaje del total dentro de cada grupo}

Una pregunta clásica: ``¿qué porcentaje representa cada cosa \emph{dentro}
de su grupo?''. La receta es: primero resumes con \codigo{summarise()} y
después, con el resultado todavía agrupado por el nivel superior, calculas
el porcentaje con \codigo{mutate()}.

\begin{ejemplo}{Mezcla de métodos de pago por trimestre}
Finanzas negocia las comisiones bancarias y quiere saber qué porcentaje de
los ingresos de cada trimestre entró por cada método de pago.
\end{ejemplo}

\begin{rcode}
mezcla_pago <- transacciones %>%
  group_by(trimestre, metodo_pago) %>%
  summarise(ingresos = sum(total_transaccion),
            .groups = "drop_last") %>%   # queda agrupado por trimestre
  mutate(pct_trimestre = round(ingresos / sum(ingresos) * 100, 1)) %>%
  ungroup()                               # quitamos el último grupo
mezcla_pago %>% filter(trimestre %in% c("Q1", "Q4"))
\end{rcode}
\begin{rsalida}
# A tibble: 8 x 4
  trimestre metodo_pago     ingresos pct_trimestre
  <chr>     <chr>              <dbl>         <dbl>
1 Q1        Efectivo          80565.          33.7
2 Q1        Tarjeta Crédito   88269.          36.9
3 Q1        Tarjeta Débito    51032.          21.3
4 Q1        Transferencia     19524.           8.2
5 Q4        Efectivo          78777.          31.9
6 Q4        Tarjeta Crédito   87412.          35.4
7 Q4        Tarjeta Débito    59178.          24  
8 Q4        Transferencia     21537.           8.7
\end{rsalida}

Con \codigo{.groups = "drop\_last"} el resultado queda agrupado por
\codigo{trimestre}, así que \codigo{sum(ingresos)} dentro de
\codigo{mutate()} es la suma \emph{de ese trimestre} y los porcentajes de
cada trimestre suman 100. La tarjeta de débito ganó terreno: pasó de
21.3\% de los ingresos en el Q1 a 24\% en el Q4, mientras que el efectivo
bajó de 33.7\% a 31.9\%. Si hubieras quitado los grupos antes del \codigo{mutate()},
\codigo{sum(ingresos)} sería el total del año y los porcentajes serían
respecto al año completo: ambas preguntas son válidas, pero distintas.

\subsubsection{Agrupar sin resumir: \texorpdfstring{\codigo{group\_by()}}{group\_by()} + \texorpdfstring{\codigo{mutate()}}{mutate()}}

\codigo{group\_by()} también funciona con \codigo{mutate()}: en lugar de
reducir cada grupo a una fila, conserva todas las filas y calcula dentro de
cada grupo. Es perfecto para comparar cada registro contra el promedio de su
grupo.

\begin{ejemplo}{¿Qué ventas están muy por encima del promedio de su tienda?}
Para detectar ventas atípicas no basta con compararlas contra el promedio
general: una venta de \$3,000 puede ser normal en una tienda y excepcional
en otra.
\end{ejemplo}

\begin{rcode}
ventas_vs_tienda <- ventas %>%
  group_by(tienda_id) %>%
  mutate(
    promedio_tienda = mean(total),              # promedio de SU tienda
    veces_promedio  = round(total / promedio_tienda, 2)
  ) %>%
  ungroup()                                     # ¡no lo olvides!

ventas_vs_tienda %>%
  select(fecha, tienda_id, total, promedio_tienda, veces_promedio) %>%
  slice_max(veces_promedio, n = 4)
\end{rcode}
\begin{rsalida}
# A tibble: 4 x 5
  fecha      tienda_id total promedio_tienda veces_promedio
  <date>         <dbl> <dbl>           <dbl>          <dbl>
1 2023-04-21         2 4925.           1443.           3.41
2 2023-07-24         8 4950.           1493.           3.32
3 2023-02-16         9 4367.           1334.           3.27
4 2023-02-15         8 4638.           1493.           3.11
\end{rsalida}

La venta del 21 de abril en la tienda 2 fue 3.41 veces el promedio de
esa tienda: vale la pena confirmar que sea real. Observa que el resultado
conserva las 365 filas; \codigo{mutate()} no resume, solo agrega columnas.

\begin{advertencia}
Una tabla agrupada \emph{recuerda} sus grupos. Si olvidas
\codigo{ungroup()}, los pasos siguientes (un \codigo{mutate()}, un
\codigo{slice\_max()}, un \codigo{summarise()}) se aplicarán grupo por grupo
sin que te des cuenta. Cuando imprimas una tabla y veas una línea
\texttt{\# Groups: tienda\_id [10]} debajo del encabezado, esa tabla sigue
agrupada.
\end{advertencia}

\subsubsection{Valores faltantes en los resúmenes}

Si una columna tiene aunque sea un \codigo{NA}, su suma y su promedio serán
\codigo{NA}: R se niega a adivinar. Usa \codigo{na.rm = TRUE} para ignorar
los faltantes, pero hazlo \emph{a conciencia} y reporta cuántos había.

\begin{rcode}
metas %>%
  summarise(
    suma_sin_na_rm  = sum(meta),               # NA "contagia"
    suma_capturada  = sum(meta, na.rm = TRUE), # ignora los NA
    metas_faltantes = sum(is.na(meta))         # cuántos NA hay
  )
\end{rcode}
\begin{rsalida}
# A tibble: 1 x 3
  suma_sin_na_rm suma_capturada metas_faltantes
           <dbl>          <dbl>           <int>
1             NA         153000               2
\end{rsalida}

El truco \codigo{sum(is.na(x))} cuenta los faltantes: \codigo{is.na()}
devuelve \codigo{TRUE}/\codigo{FALSE} y al sumar, cada \codigo{TRUE} vale 1.

\begin{hazlotu}
Con \codigo{transacciones}: (a) calcula por \codigo{dia\_semana} el número
de transacciones, los ingresos y el ticket promedio, ordenado de mayor a
menor ingreso; (b) con \codigo{count()} y \codigo{wt}, obtén los ingresos
por tienda; (c) calcula qué porcentaje de las transacciones de cada tienda
se pagó en efectivo.
\end{hazlotu}

% ----------------------------------------------------------------------------
\section{Funciones de ventana para KPIs}
\label{sec:m2-ventana}

Las funciones de resumen (\codigo{sum()}, \codigo{mean()}\ldots) convierten
muchos valores en uno. Las \termino{funciones de ventana} (\emph{window
functions}) devuelven \emph{un valor por cada fila}, pero calculado mirando
a otras filas: la anterior, la siguiente, todas las previas o todo el grupo.
Con ellas se calculan casi todos los KPIs de tendencia de un tablero de BI.

\begin{table}[H]
\centering
\caption{Funciones de ventana más útiles en BI.}
\label{tab:m2-ventana}
\small
\begin{tabularx}{\textwidth}{@{}lXl@{}}
\toprule
Función & Qué devuelve & KPI típico \\
\midrule
\codigo{lag(x)} & El valor de la fila anterior & Crecimiento mes contra mes \\
\codigo{lead(x)} & El valor de la fila siguiente & Comparar con el periodo siguiente \\
\codigo{cumsum(x)} & Suma acumulada & Acumulado del año (YTD) \\
\codigo{row\_number()} & Posición 1, 2, 3\ldots{} sin empates & Numerar filas \\
\codigo{min\_rank()} & Ranking con huecos tras empates (1, 2, 2, 4) & Ranking de tiendas \\
\codigo{dense\_rank()} & Ranking sin huecos (1, 2, 2, 3) & Niveles de precio \\
\codigo{x / sum(x)} & Proporción de cada fila en el total & Participación de mercado \\
\bottomrule
\end{tabularx}
\end{table}

Primero construimos la serie mensual de ingresos: una fila por mes.

\begin{rcode}
ventas_mensuales <- ventas %>%
  group_by(mes) %>%
  summarise(ingresos = sum(total), num_ventas = n(),
            .groups = "drop") %>%
  arrange(mes)                  # ¡las ventanas dependen del orden!
\end{rcode}

\subsection{Crecimiento mes contra mes (MoM) y acumulado del año (YTD)}

El \termino{crecimiento mes contra mes} (MoM, \emph{month over month}) mide
cuánto cambió un indicador respecto al mes anterior:
\[
\text{MoM} = \left(\frac{\text{ingresos del mes}}{\text{ingresos del mes
anterior}} - 1\right) \times 100 .
\]
El \termino{acumulado del año} (YTD, \emph{year to date}) es la suma de
todos los meses del año hasta el mes actual.

\begin{rcode}
kpis_mensuales <- ventas_mensuales %>%
  mutate(
    mes_anterior = lag(ingresos),                    # fila previa
    mom_pct      = (ingresos / lag(ingresos) - 1) * 100,
    ytd          = cumsum(ingresos)                  # acumulado
  )

kpis_mensuales %>%
  mutate(across(where(is.numeric), ~ round(.x, 1))) %>%
  select(mes, ingresos, mes_anterior, mom_pct, ytd)
\end{rcode}
\begin{rsalida}
# A tibble: 12 x 5
   mes     ingresos mes_anterior mom_pct     ytd
   <chr>      <dbl>        <dbl>   <dbl>   <dbl>
 1 2023-01   46884.          NA     NA    46884.
 2 2023-02   42656.       46884.    -9    89540.
 3 2023-03   36734.       42656.   -13.9 126274.
 4 2023-04   34738.       36734.    -5.4 161012.
 5 2023-05   48108.       34738.    38.5 209120.
 6 2023-06   37871.       48108.   -21.3 246991.
 7 2023-07   54056.       37871.    42.7 301047.
 8 2023-08   41421        54056.   -23.4 342468.
 9 2023-09   38950.       41421     -6   381418.
10 2023-10   36584.       38950.    -6.1 418002.
11 2023-11   50398.       36584.    37.8 468400.
12 2023-12   45516.       50398.    -9.7 513917.
\end{rsalida}

¿Qué nos dice esto? Enero no tiene mes anterior dentro de los datos, por eso
\codigo{lag()} devuelve \codigo{NA} y el crecimiento también: es correcto,
no un error. El mayor crecimiento fue en julio (42.7\% contra junio) y la
peor caída en agosto ($-23.4$\%). Nota el patrón de ``serrucho'': a cada
mes fuerte le sigue uno débil. La columna \codigo{ytd} muestra que al cierre
de junio se llevaban \$246,991 (48\% del año) y el año cerró en
\$513,916.60.

\begin{advertencia}
\codigo{lag()} toma ``la fila de arriba'', no ``el mes anterior''. Si la
tabla no está ordenada por mes, o si falta un mes, el cálculo será
incorrecto sin ningún aviso. Por eso ordenamos con \codigo{arrange(mes)}
antes, y en la sección~\ref{sec:m2-complete} verás cómo rellenar los meses
faltantes con \codigo{complete()}. Para una serie con fechas desordenadas
puedes usar \codigo{lag(ingresos, order\_by = mes)}.
\end{advertencia}

\subsection{Participación y promedio móvil}

La \termino{participación} de cada mes es su porcentaje del total anual. El
\termino{promedio móvil} de 3 meses suaviza los altibajos: en cada mes
promedia ese mes y los dos anteriores, y deja ver la tendencia de fondo.
\codigo{lag(x, 2)} es el valor de dos filas atrás.

\begin{rcode}
kpis_mensuales <- kpis_mensuales %>%
  mutate(
    participacion = ingresos / sum(ingresos) * 100,
    prom_movil_3m = (ingresos + lag(ingresos) + lag(ingresos, 2)) / 3,
    vs_siguiente  = lead(ingresos) - ingresos     # ¿el mes que sigue?
  )

kpis_mensuales %>%
  mutate(across(where(is.numeric), ~ round(.x, 1))) %>%
  select(mes, ingresos, participacion, prom_movil_3m, vs_siguiente)
\end{rcode}
\begin{rsalida}
# A tibble: 12 x 5
   mes     ingresos participacion prom_movil_3m vs_siguiente
   <chr>      <dbl>         <dbl>         <dbl>        <dbl>
 1 2023-01   46884.           9.1           NA        -4228.
 2 2023-02   42656.           8.3           NA        -5923.
 3 2023-03   36734.           7.1        42091.       -1996.
 4 2023-04   34738.           6.8        38042.       13370.
 5 2023-05   48108.           9.4        39860.      -10238.
 6 2023-06   37871.           7.4        40239        16185.
 7 2023-07   54056.          10.5        46678.      -12635.
 8 2023-08   41421            8.1        44449.       -2471.
 9 2023-09   38950.           7.6        44809.       -2366 
10 2023-10   36584.           7.1        38985.       13814.
11 2023-11   50398.           9.8        41977.       -4882.
12 2023-12   45516.           8.9        44166.          NA 
\end{rsalida}

Cada mes aporta entre 6.8\% (abril) y 10.5\% (julio) de las ventas del
año. El promedio móvil empieza en marzo (necesita tres meses) y se mueve
entre \$38,042 y \$46,678: un rango mucho más estrecho que el de los
ingresos mensuales, que van de \$34,738 a \$54,056. Es decir, detrás de la volatilidad mensual el negocio
es bastante estable. La columna \codigo{vs\_siguiente}, calculada con
\codigo{lead()}, es \codigo{NA} en diciembre porque no hay mes siguiente.

\subsection{Rankings}

Hay tres funciones de ranking y se diferencian en cómo tratan los empates.
Míralo con un ejemplo pequeño:

\begin{rcode}
tibble(vendedor = c("Ana", "Beto", "Caro", "Dani", "Eva"),
       ventas   = c(900, 700, 700, 500, 300)) %>%
  mutate(
    row_number = row_number(desc(ventas)),  # 1,2,3,4,5 (sin empates)
    min_rank   = min_rank(desc(ventas)),    # 1,2,2,4,5 (con hueco)
    dense_rank = dense_rank(desc(ventas))   # 1,2,2,3,4 (sin hueco)
  )
\end{rcode}
\begin{rsalida}
# A tibble: 5 x 5
  vendedor ventas row_number min_rank dense_rank
  <chr>     <dbl>      <int>    <int>      <int>
1 Ana         900          1        1          1
2 Beto        700          2        2          2
3 Caro        700          3        2          2
4 Dani        500          4        4          3
5 Eva         300          5        5          4
\end{rsalida}

Usamos \codigo{desc()} para que el valor más alto reciba el lugar 1.
\codigo{min\_rank()} es el ranking ``deportivo'' (dos segundos lugares y
nadie en tercero) y es el más usado en reportes.

\begin{ejemplo}{Las 3 mejores tiendas de cada trimestre}
Recursos Humanos quiere premiar cada trimestre a las tres tiendas con más
ingresos. ¿Cuáles son y repiten las mismas?
\end{ejemplo}

\begin{rcode}
podio_trimestral <- ventas_trim_tienda %>%     # la calculamos antes
  group_by(trimestre) %>%
  mutate(lugar = min_rank(desc(ingresos))) %>% # ranking por trimestre
  filter(lugar <= 3) %>%
  arrange(trimestre, lugar) %>%
  ungroup()
podio_trimestral
\end{rcode}
\begin{rsalida}
# A tibble: 12 x 4
   trimestre tienda_id ingresos lugar
   <chr>         <dbl>    <dbl> <int>
 1 Q1                8   17226.     1
 2 Q1                1   16711.     2
 3 Q1                9   16706.     3
 4 Q2                2   20947.     1
 5 Q2                8   19290.     2
 6 Q2                3   14997.     3
 7 Q3                9   19581.     1
 8 Q3                6   18015.     2
 9 Q3                2   16639.     3
10 Q4                5   19978.     1
11 Q4                8   18245.     2
12 Q4                1   14823.     3
\end{rsalida}

La tienda 8 subió al podio en tres de los cuatro trimestres: es la más
consistente. Más interesante es la tienda 1: en el ranking anual quedó en
octavo lugar, pero estuvo en el podio del Q1 y del Q4. Sus problemas se
concentran a mitad de año, algo que el total anual escondía. Observa también
que \codigo{min\_rank()} dentro de \codigo{group\_by(trimestre)} reinicia el
ranking en cada trimestre.

\begin{ejemplo}{Análisis ABC (Pareto) de productos}
El principio de Pareto dice que pocos productos generan la mayor parte de las
ventas. El \termino{análisis ABC} clasifica los productos en: \textbf{A}, los
que acumulan el primer 80\% de los ingresos; \textbf{B}, el siguiente 15\%;
y \textbf{C}, el último 5\%. ¿Cuántos productos hay en cada clase?
\end{ejemplo}

\begin{rcode}
abc_productos <- ventas %>%
  group_by(producto_id) %>%
  summarise(ingresos = sum(total), .groups = "drop") %>%
  arrange(desc(ingresos)) %>%                    # del mejor al peor
  mutate(
    pct_acumulado = cumsum(ingresos) / sum(ingresos) * 100,
    clase = case_when(
      pct_acumulado <= 80 ~ "A",
      pct_acumulado <= 95 ~ "B",
      .default            = "C"
    )
  )

abc_productos %>%
  group_by(clase) %>%
  summarise(productos = n(), ingresos = sum(ingresos),
            .groups = "drop") %>%
  mutate(pct_ingresos = round(ingresos / sum(ingresos) * 100, 1))
\end{rcode}
\begin{rsalida}
# A tibble: 3 x 4
  clase productos ingresos pct_ingresos
  <chr>     <int>    <dbl>        <dbl>
1 A            31  402751.         78.4
2 B            12   83751.         16.3
3 C             7   27415.          5.3
\end{rsalida}

Aquí \codigo{cumsum()} se aplica después de ordenar de mayor a menor, así que
\codigo{pct\_acumulado} dice qué porcentaje de los ingresos suman los
mejores productos hasta esa fila. El resultado es interesante porque
\emph{no} sigue el 80/20 clásico: se necesitan 31 de los 50 productos (62\%)
para juntar el 78.4\% de los ingresos. Las ventas de esta empresa están
poco concentradas, así que descontinuar productos de clase C (7 productos,
5.3\% de los ingresos) tendría poco impacto en la venta. En una empresa real
este análisis guía decisiones de inventario: a los productos A nunca les
debe faltar existencia.

\begin{hazlotu}
Repite el cálculo de MoM y YTD, pero con \codigo{transacciones} y la columna
\codigo{total\_transaccion}. ¿Cuál fue el mejor mes? ¿En qué mes el
acumulado superó el medio millón de pesos? Pista: después de calcular
\codigo{ytd}, usa \codigo{filter(ytd > 500000) \%>\% slice\_head(n = 1)}.
\end{hazlotu}
